\section{Related Works}

%% use \paragraph{...}

\paragraph{On-Device LLM Inference with NPUs.} \texttt{llm.npu}~\cite{mllm-npu} pioneered the use of per-tensor quantized INT8 GEMM on NPUs to accelerate the prefill phase of LLMs, employing the CPU to assist in outlier-related computations to maintain accuracy. HeteroLLM~\cite{heterollm} achieves collaborative inference between GPUs and NPUs through tensor partitioning. PowerServe~\cite{powerserve}, an open-source inference framework, leverages the intermediate ONNX format to implement custom quantized and floating-point computation partitioning. ShadowAttn~\cite{shadowattn} utilizes collaboration between NPUs and CPUs/GPUs to accelerate sparse attention. ExecuTorch~\cite{executorch,torchao} is a well-known open-source edge-side DNN inference framework, supporting SpinQuant~\cite{spinquant} and the QNN backend. All of the above works are based on Qualcomm’s Hexagon NPU and use the closed-source QNN~\cite{qualcomm-qnn} as the backend. Works from vivo~\cite{bluelm,edgeinfinite} have utilized MediaTek’s NPU, but due to the non-public accessibility of the NeuroPilot SDK, such research remains scarce.

\paragraph{LLM Quantization.} The most well-known post-training quantization algorithms include GPTQ~\cite{gptq} and AWQ~\cite{awq}, they perform weight-only quantization and only require small amounts of calibration data, therefore they are extensively used. Subsequent methods such as SmoothQuant~\cite{smoothquant}, DuQuant~\cite{duquant} broadened the scope to include weight-activation quantization, tackling the more challenging task of quantizing activations by developing techniques to mitigate their problematic outlier distributions. Most recently, comprehensive approaches such as QuaRot~\cite{quarot} and SpinQuant~\cite{spinquant} have emerged, aiming to quantize all major components - weights, activations, and the critical KV cache - often down to 4-bits. These methods leverage rotation transformations to create more quantization-friendly feature distributions throughout the model.

\paragraph{Speculative Decoding.} Speculative Decoding~\cite{spec-decode,specinfer,zengyixiao-sd} is a class of acceleration methods for LLM inference, the core of which is to verify multiple speculated tokens in one model forward pass to alleviate the memory-bound issue of LLM decoding. There are various extended variants of Speculative Decoding, and some~\cite{hao2024hybrid,reward-guided-sd} no longer strictly follow the distribution of the target model. In theory, generalized Speculative Decoding and test-time scaling methods both belong to the generalized Generate-then-Verify framework, and our system can theoretically support these applications seamlessly.
